\documentclass[12pt,a4paper]{article}

\usepackage{amsmath,amssymb,amsthm}
\usepackage{geometry}
\usepackage{booktabs}
\usepackage{array}
\usepackage{hyperref}
\usepackage{setspace}
\usepackage{parskip}

\geometry{margin=1in}
\onehalfspacing

\newtheorem{proposition}{Proposition}
\newtheorem{corollary}{Corollary}
\newtheorem{definition}{Definition}
\newtheorem{remark}{Remark}

\title{\textbf{Wage Inequality, Upskilling, and Long-Run Growth:\\
A Three-Player Model with Linearly Approximated Risk Premia}\\[0.8em]
\large Working Paper}
\author{}
\date{\today}

\begin{document}

\maketitle

\begin{abstract}
We develop a tractable three-player model in which a single wage inequality
parameter $\mu \in (0,\tfrac{1}{2})$ --- the income share accruing to the
lower-ranked worker group --- governs both worker behaviour and firm asset
accumulation. Workers in two income groups allocate wages between upskilling
investment and participation in a zero-sum displacement lottery. Risk
preferences are captured through a linearly approximated Friedman-Savage
utility, which preserves the key asymmetry that lower-income workers are
risk-seeking and higher-income workers are risk-averse, while keeping all
equilibrium objects in closed form. A representative firm chooses $\mu$ each
period to maximise asset accumulation subject to a restructuring cost. We
address two questions. First, how do workers respond to wage inequality $\mu$:
specifically, how do upskilling investment, lottery participation, and
risk premia change as $\mu$ varies? Second, what constraints does $\mu$ impose
on long-term growth: what is the feasible range of $\mu$, what is the optimal
$\mu^*$ chosen by the firm, and when does the system admit a positive
steady-state asset level versus collapsing to a poverty trap? We derive
fully explicit steady-state equilibria for all three players and a complete
set of signed comparative statics.
\end{abstract}

\tableofcontents
\newpage

%----------------------------------------------------------------------
\section{Introduction}
%----------------------------------------------------------------------

Wage inequality is typically treated in macroeconomic models as either an
exogenous parameter or an outcome of competitive factor markets. In either
case it is rarely connected to the investment decisions of workers or to the
strategic incentives of the firm setting wages. This paper proposes a unified
framework in which wage inequality --- captured by a single parameter $\mu$ ---
is the central object connecting all three players in the economy: a lower-income
worker, a higher-income worker, and a firm.

The model has two core mechanisms. The first is a worker-level allocation
problem. Each worker receives income determined by $\mu$ and divides it between
upskilling investment, which raises future productive capacity, and participation
in a zero-sum displacement lottery, which redistributes income between the two
workers. The key friction is a participation constraint: each worker must invest
enough in upskilling to maintain a subsistence base income. This constraint binds
at equilibrium and determines upskilling investment as a direct function of $\mu$.

The second mechanism is a firm-level asset accumulation problem. The firm
accumulates productive assets and pays wages as a share of those assets.
It chooses $\mu$ each period to maximise next-period assets, internalising
how $\mu$ affects the workers' upskilling output, which in turn feeds back into
asset growth. The firm faces a convex restructuring cost for deviating from a
reference wage split $\mu_0$, creating an interior optimum.

The two research questions addressed are as follows. The first asks how workers
respond to wage inequality. A change in $\mu$ alters the income of both workers
simultaneously. For worker $L$ it changes the base income available for upskilling
and the lottery stake available for the displacement mechanism. For worker $H$ the
same logic applies in the opposite direction. The model delivers sharp predictions
about the direction and magnitude of these responses, including a non-obvious
result: total upskilling investment $K^* = 2\kappa^{\min}$ is invariant to $\mu$
at the constrained equilibrium, even though the fractions $\lambda_L^*$ and
$\lambda_H^*$ move in opposite directions. Redistribution shifts the composition
of upskilling but not its aggregate level.

The second question asks what constraints $\mu$ places on long-term growth. Not
all values of $\mu$ support a positive steady-state asset level. A viability
condition must hold, and $\mu$ enters this condition through the upskilling index
$\Phi(\mu)$. The firm's optimal $\mu^*$ is increasing in the asset stock --- richer
firms enforce greater equality --- and the steady-state gap $\mu^* - \mu_0$ is
determined entirely by primitives. Below the viability threshold, no positive
steady state exists and the firm's asset stock collapses regardless of initial
conditions, constituting a poverty trap.

The model uses a linear approximation of the second derivative of the
Friedman-Savage utility function. This is the middle ground between a fully
linear (risk-neutral) utility, which eliminates all preference asymmetry between
workers, and the full cubic Friedman-Savage specification, which preserves
exact curvature but does not admit closed-form risk premia. Under the linear
approximation, risk premia remain explicit and retain the correct sign asymmetry:
the lower-income worker has a negative risk premium (risk-seeking) and the
higher-income worker has a positive risk premium (risk-averse), with the
crossover occurring exactly when subsistence equals the Friedman-Savage
inflection point.

%----------------------------------------------------------------------
\section{Model Setup}
%----------------------------------------------------------------------

\subsection{Players, Incomes, and the Role of $\mu$}

There are three players: worker $L$, worker $H$, and a firm. The firm holds
asset stock $A_t$ and pays a total wage bill $\pi A_t$ each period, where
$\pi \in (0,1)$ is the wage share. Let $R$ denote total income available to
workers, combining a base component and the wage bill. The firm chooses the
income-split parameter $\mu \in (0, \tfrac{1}{2})$, which allocates income as:
\begin{equation}
  M_L = R\mu, \qquad M_H = R(1-\mu)
  \label{eq:incomes}
\end{equation}
Since $\mu < \tfrac{1}{2}$, worker $H$ always earns more than worker $L$.
The parameter $\mu$ is the single measure of wage inequality in the model:
lower $\mu$ means greater inequality, higher $\mu$ means greater equality.
The firm treats $\mu$ as a policy instrument, choosing it optimally each period.

\begin{definition}[Wage inequality parameter]
$\mu \in (0, \tfrac{1}{2})$ is the income share of worker $L$. The income
ratio is $M_H/M_L = (1-\mu)/\mu > 1$, which is strictly decreasing in $\mu$.
Greater wage equality corresponds to higher $\mu$.
\end{definition}

The parameter $\mu$ has a microfoundation. In the underlying $N$-worker
displacement model with a Pareto income ladder $m_i = x_0(1-i/(N+1))^{-1/\alpha}$
and a Friedman-Savage class barrier at rank $i^*$, the two-player reduction
implies:
\begin{equation}
  \mu \approx \frac{\int_0^{i^*/N} x_0(1-s)^{-1/\alpha}\,ds}
                   {\int_0^{1} x_0(1-s)^{-1/\alpha}\,ds}
  \label{eq:mu_micro}
\end{equation}
so that $\mu$ is the lower hump's share of total representative income.
In the macroeconomic model, $\mu$ is taken as the firm's choice variable
rather than a structural outcome, allowing us to study optimal policy.

\subsection{Worker Budget and Allocation}

Each worker $s \in \{L, H\}$ receives income $M_s$ and allocates fraction
$\lambda_s \in [0,1]$ to upskilling and the remainder to the lottery:
\begin{equation}
  \kappa_s = \lambda_s M_s \quad \text{(upskilling investment)},
  \qquad
  \ell_s = (1-\lambda_s)M_s \quad \text{(lottery stake)}
  \label{eq:budget}
\end{equation}
Upskilling produces output $f(\kappa_s) = \kappa_s^\gamma$ with $\gamma \in (0,1)$
(diminishing returns). The certain base income after upskilling and human capital
depreciation $\delta_w > 0$ is:
\begin{equation}
  b_s(\lambda_s) = \kappa_s + \kappa_s^\gamma - \delta_w
  \label{eq:base_income}
\end{equation}

\subsection{Zero-Sum Displacement Lottery}

The lottery represents a displacement mechanism inherited from the underlying
$N$-worker stochastic model. Workers contribute stakes to a common pool, and
one worker wins the entire pool each period. The pool is:
\begin{equation}
  \Lambda = W(1-\lambda_L)M_L + (1-\lambda_H)M_H = W\ell_L + \ell_H
  \label{eq:pool}
\end{equation}
where $W \geq 1$ is a population weight equal to $i^*/(N-i^*)$ in the
$N$-worker foundation, reflecting that the lower-ranked group is more numerous
and generates more displacement events per period. Each worker wins the pool
with probability $\tfrac{1}{2}$. The outcome incomes are:
\begin{align}
  M_L^{\mathrm{win}}  &= b_L + (1-\lambda_H)M_H \label{eq:MLwin} \\
  M_L^{\mathrm{lose}} &= b_L - W(1-\lambda_L)M_L \label{eq:MLlose} \\
  M_H^{\mathrm{win}}  &= b_H + W(1-\lambda_L)M_L \label{eq:MHwin} \\
  M_H^{\mathrm{lose}} &= b_H - (1-\lambda_H)M_H \label{eq:MHlose}
\end{align}
The lottery is zero-sum: whatever worker $L$ gains, worker $H$ loses, and
vice versa. Total income is conserved by the lottery and only upskilling
changes the aggregate.

\subsection{Participation Constraint}

Each worker must maintain a base income above a subsistence floor
$\bar{b}(\mu) = \bar{b}_0 M_L = \bar{b}_0 R\mu$, which is proportional
to the lower-income worker's wage. This ensures the floor is endogenous
to wage inequality: as $\mu$ rises, the subsistence requirement rises,
forcing more upskilling investment. The participation constraint is:
\begin{equation}
  b_s(\lambda_s) \geq \bar{b}(\mu) \iff \kappa_s + \kappa_s^\gamma \geq \bar{b}_0 R\mu + \delta_w
  \label{eq:pc}
\end{equation}
The minimum upskilling investment satisfying this constraint, $\kappa^{\min}(\mu)$,
solves:
\begin{equation}
  \kappa^{\min} + \left(\kappa^{\min}\right)^\gamma = \bar{b}_0 R\mu + \delta_w
  \label{eq:kappa_min}
\end{equation}
For $\gamma = 2$, which we adopt throughout for closed-form solutions:
\begin{equation}
  \kappa^{\min}(\mu) = \frac{-1 + \sqrt{1 + 4(\bar{b}_0 R\mu + \delta_w)}}{2}
  \label{eq:kappa_min_gamma2}
\end{equation}
This is strictly increasing and concave in $\mu$. The sensitivity of the
upskilling floor to wage inequality is:
\begin{equation}
  \frac{d\kappa^{\min}}{d\mu} = \frac{\bar{b}_0 R}{1 + 2\kappa^{\min}(\mu)} \equiv \Psi(\mu) > 0
  \label{eq:dkdmu}
\end{equation}

%----------------------------------------------------------------------
\section{Worker Preferences and the Linear Risk Premium Approximation}
%----------------------------------------------------------------------

\subsection{Friedman-Savage Utility}

Worker utility is defined over next-period income. We use a cubic
Friedman-Savage specification with a single inflection point $m^*$:
\begin{equation}
  u(m) = m - \frac{\alpha_u}{2}(m-m^*)^2 + \frac{\beta_u}{3}(m-m^*)^3
  \label{eq:cubic_u}
\end{equation}
where $\alpha_u, \beta_u > 0$. The key properties are:
\begin{itemize}
  \item $u''(m) = -\alpha_u + 2\beta_u(m-m^*)$, which is negative for
        $m > m^*$ (concave, risk-averse) and positive for $m < m^*$
        (convex, risk-seeking).
  \item The inflection at $m^*$ is the income level at which risk attitude
        switches. Worker $L$ with $M_L = R\mu < m^*$ is risk-seeking;
        worker $H$ with $M_H = R(1-\mu) > m^*$ is risk-averse.
\end{itemize}

\subsection{The Linear Approximation of $u''$}

Rather than using the full cubic $u''$, we linearise the second derivative
around the subsistence income $\bar{b}$:
\begin{equation}
  u''(m) \approx u''(\bar{b}) + u'''(\bar{b})\cdot(m - \bar{b})
  = \underbrace{(-\alpha_u + 2\beta_u d)}_{\text{curvature at }\bar{b}}
  + \underbrace{2\beta_u}_{\text{slope}} \cdot (m - \bar{b})
  \label{eq:linear_upp}
\end{equation}
where $d = \bar{b} - m^*$ is the deviation of subsistence from the inflection.
Since $u'''(m) = 2\beta_u$ is already constant under the cubic specification,
this linearisation is exact for the cubic --- it introduces no approximation
error for $u''$ itself. What it approximates is the Arrow-Pratt coefficient
$A(m) = -u''(m)/u'(m)$, where we replace the denominator $u'(m)$ with $u'(\bar{b})$:
\begin{equation}
  A(m) \approx \tilde{A}(m) \equiv \frac{-u''(m)}{u'(\bar{b})}
  = \frac{\alpha_u - 2\beta_u(m-m^*)}{u'(\bar{b})}
  \label{eq:AP_linear}
\end{equation}
This approximation is valid when $u'$ does not vary much over the lottery
outcome range, which holds when prizes are small relative to the distance
from $\bar{b}$ to $m^*$.

\subsection{Linearly Approximated Risk Premia}

For a binary lottery paying $\bar{b} + \Pi_{-s}$ or $\bar{b} - W_s\ell_s$
with equal probability, the exact risk premium is:
\begin{equation}
  \rho_s = -\frac{\Pi_s^2}{8} \cdot A_s(\bar{b} + \epsilon_s)
  \label{eq:rho_exact}
\end{equation}
where $\Pi_s = \Pi_L = \Pi_H = \Lambda^* = Wa + c$ is the total lottery
spread (equal for both workers since both win or lose the entire pool),
and $\epsilon_s$ is the net expected transfer. Under the linear approximation
(\ref{eq:AP_linear}), the risk premium becomes:
\begin{equation}
  \tilde{\rho}_s = -\frac{\Pi^2}{8} \cdot \tilde{A}(\bar{b} + \epsilon_s)
  = \frac{\Pi^2}{8} \cdot \frac{\alpha_u - 2\beta_u D_s}{u'(\bar{b})}
  \label{eq:rho_linear}
\end{equation}
where $D_s \equiv \bar{b} - m^* + \epsilon_s$ is the deviation of worker $s$'s
expected lottery outcome from the inflection point, and we use the
compressed notation:
\begin{align}
  a &= M_L - \kappa^{\min} = R\mu - \kappa^{\min}(\mu) \quad \text{(worker $L$ lottery stake)}
  \label{eq:a} \\
  c &= M_H - \kappa^{\min} = R(1-\mu) - \kappa^{\min}(\mu) \quad \text{(worker $H$ lottery stake)}
  \label{eq:c} \\
  \epsilon_L &= \tfrac{c - Wa}{2} \quad \text{(net expected lottery transfer to $L$)}
  \label{eq:epsilonL} \\
  \epsilon_H &= \tfrac{Wa - c}{2} = -\epsilon_L \quad \text{(net expected lottery transfer to $H$)}
  \label{eq:epsilonH}
\end{align}
The sign of the risk premia is determined by $\alpha_u - 2\beta_u D_s$. Since
$D_L = \bar{b} - m^* + \epsilon_L < 0$ (worker $L$ below inflection) and
$D_H = \bar{b} - m^* + \epsilon_H > 0$ (worker $H$ above inflection), we have:
\begin{align}
  \tilde{\rho}_L &< 0 \quad \text{(worker $L$ risk-seeking: pays to enter lottery)}
  \label{eq:rhoL_sign} \\
  \tilde{\rho}_H &> 0 \quad \text{(worker $H$ risk-averse: demands compensation)}
  \label{eq:rhoH_sign}
\end{align}
The sign asymmetry between the two workers is preserved entirely. The net
social risk premium under the linear approximation is:
\begin{equation}
  \tilde{\rho}_L + \tilde{\rho}_H
  = \frac{\Pi^2}{4} \cdot \frac{u''(\bar{b})}{u'(\bar{b})}
  = -\frac{\Pi^2}{4} \cdot A(\bar{b})
  \label{eq:social_premium}
\end{equation}
This collapses to a single expression: the net social premium equals
$-\Pi^2/4$ times the Arrow-Pratt coefficient at subsistence. The lottery is
socially desirable (net positive premium to participants) if and only if
$A(\bar{b}) < 0$, i.e.\ $\bar{b} < m^*$. The lottery mechanism generates
aggregate welfare gains precisely when the representative subsistence level
is below the Friedman-Savage inflection --- a clean and intuitive condition.

\subsection{Equilibrium Utility Under Linear Approximation}

The worker's expected utility objective, expanded to second order with the
linear approximation of $u''$, is:
\begin{equation}
  V_s(\lambda_s;\lambda_{-s}) \approx u(\bar{b})
  + \epsilon_s \cdot u'(\bar{b})
  + \tilde{\rho}_s
  \label{eq:V_approx}
\end{equation}
The three terms have clear interpretations. The first, $u(\bar{b})$, is the
utility of subsistence income --- the floor guaranteed by the participation
constraint. The second, $\epsilon_s \cdot u'(\bar{b})$, is the expected lottery
transfer valued at the marginal utility of subsistence income. The third,
$\tilde{\rho}_s$, is the linearly approximated risk premium, which captures the
curvature correction: worker $L$ gains from the risk ($\tilde\rho_L < 0$ means
they receive a premium for bearing it) and worker $H$ loses ($\tilde\rho_H > 0$
means they pay a premium to avoid it).

%----------------------------------------------------------------------
\section{Question 1: How Do Workers Respond to Wage Inequality $\mu$?}
%----------------------------------------------------------------------

\subsection{Optimal Worker Strategies}

Since the participation constraint $b_s \geq \bar{b}$ binds at equilibrium
(workers optimally set base income exactly to subsistence, allocating the
maximum feasible fraction to the lottery), the optimal upskilling fractions are:
\begin{equation}
  \lambda_L^*(\mu) = \frac{\kappa^{\min}(\mu)}{R\mu},
  \qquad
  \lambda_H^*(\mu) = \frac{\kappa^{\min}(\mu)}{R(1-\mu)}
  \label{eq:lambda_star}
\end{equation}
These are fully explicit functions of $\mu$ through $\kappa^{\min}(\mu)$ in
equation (\ref{eq:kappa_min_gamma2}). Since $\mu < \tfrac{1}{2}$, it follows
that $R\mu < R(1-\mu)$ and therefore:
\begin{equation}
  \lambda_L^*(\mu) > \lambda_H^*(\mu) \quad \forall\, \mu \in \left(0,\tfrac{1}{2}\right)
  \label{eq:lambda_ordering}
\end{equation}
The lower-income worker always invests a larger fraction of income in upskilling,
not because of any intrinsic preference difference but because their lower income
means the same absolute upskilling floor $\kappa^{\min}$ represents a larger
fraction of their budget.

\subsection{Worker Responses to a Change in $\mu$}

Differentiating equations (\ref{eq:lambda_star}) with respect to $\mu$:
\begin{align}
  \frac{d\lambda_L^*}{d\mu}
  &= \frac{1}{R\mu}\frac{d\kappa^{\min}}{d\mu} - \frac{\kappa^{\min}}{R\mu^2}
  = \frac{1}{R\mu}\left[\Psi(\mu) - \frac{\kappa^{\min}(\mu)}{\mu}\right]
  \label{eq:dlL_dmu} \\[6pt]
  \frac{d\lambda_H^*}{d\mu}
  &= \frac{1}{R(1-\mu)}\frac{d\kappa^{\min}}{d\mu} + \frac{\kappa^{\min}}{R(1-\mu)^2}
  = \frac{1}{R(1-\mu)}\left[\Psi(\mu) + \frac{\kappa^{\min}(\mu)}{1-\mu}\right]
  \label{eq:dlH_dmu}
\end{align}
Since $\Psi(\mu) > 0$ and $\kappa^{\min}(\mu) > 0$, equation (\ref{eq:dlH_dmu})
is unambiguously positive: worker $H$ always increases their upskilling fraction
when $\mu$ rises. For worker $L$, the sign depends on whether the sensitivity
$\Psi(\mu) = d\kappa^{\min}/d\mu$ exceeds $\kappa^{\min}(\mu)/\mu$. Substituting
$\Psi = \bar{b}_0 R/(1+2\kappa^{\min})$, the condition $d\lambda_L^*/d\mu > 0$
reduces to:
\begin{equation}
  \bar{b}_0 R\mu > \kappa^{\min}(1+2\kappa^{\min})
  \label{eq:lL_dmu_condition}
\end{equation}
which holds when the subsistence requirement is high relative to the current
upskilling floor --- i.e.\ when the participation constraint is relatively
tight for worker $L$. In the parametric region relevant for the model
(moderate inequality, non-trivial subsistence), this condition typically holds
and $d\lambda_L^*/d\mu > 0$ as well.

\begin{proposition}[Worker upskilling responses to inequality]
At the constrained equilibrium:
\begin{enumerate}
  \item Worker $H$ always increases their upskilling fraction $\lambda_H^*$ when
        wage equality rises ($\mu$ increases): $d\lambda_H^*/d\mu > 0$ always.
  \item Worker $L$ increases their upskilling fraction $\lambda_L^*$ when
        condition (\ref{eq:lL_dmu_condition}) holds, and decreases it otherwise.
  \item The upskilling gap $\lambda_L^* - \lambda_H^*$ narrows as $\mu$ rises,
        converging to zero as $\mu \to \tfrac{1}{2}$.
\end{enumerate}
\end{proposition}

\subsection{Total Upskilling: The Aggregate Neutrality Result}

Despite the divergent individual responses, total upskilling
$K^* = \kappa_L^* + \kappa_H^* = 2\kappa^{\min}(\mu)$ moves as follows:
\begin{equation}
  \frac{dK^*}{d\mu} = 2\Psi(\mu) = \frac{2\bar{b}_0 R}{1+2\kappa^{\min}(\mu)} > 0
  \label{eq:dK_dmu}
\end{equation}
Total upskilling investment increases with wage equality. However, this increase
comes entirely from the rising subsistence floor $\bar{b}(\mu) = \bar{b}_0 R\mu$
forcing both workers to invest more --- not from any change in incentives.
The lottery pool $\Lambda^*$, by contrast, responds to $\mu$ through both
workers' stakes:
\begin{equation}
  \frac{d\Lambda^*}{d\mu} = W\frac{da}{d\mu} + \frac{dc}{d\mu}
  = (W+1)\left(R - \Psi(\mu)\right) - WR + R
  = (W+1)(R-\Psi) - (W-1)R
  \label{eq:dLam_dmu}
\end{equation}
Simplifying: $d\Lambda^*/d\mu = R(1+W) - (W+1)\Psi - (W-1)R = 2R - (W+1)\Psi$,
which is positive when total income dominates the subsistence sensitivity.

\subsection{Lottery Stakes and Risk Premia as Functions of $\mu$}

The lottery stakes at the constrained equilibrium are:
\begin{equation}
  a(\mu) = R\mu - \kappa^{\min}(\mu), \qquad c(\mu) = R(1-\mu) - \kappa^{\min}(\mu)
  \label{eq:stakes_mu}
\end{equation}
Their responses to $\mu$ are:
\begin{equation}
  \frac{da}{d\mu} = R - \Psi(\mu) > 0
  \quad \text{(worker $L$ stake rises with equality)},
  \label{eq:da_dmu}
\end{equation}
\begin{equation}
  \frac{dc}{d\mu} = -R - \Psi(\mu) < 0
  \quad \text{(worker $H$ stake falls with equality)}
  \label{eq:dc_dmu}
\end{equation}
As $\mu$ rises, worker $L$'s lottery stake increases (their income rises faster
than the upskilling floor) while worker $H$'s stake falls (their income falls
and the floor rises). The net transfer to worker $L$ is:
\begin{equation}
  \epsilon_L(\mu) = \frac{c(\mu) - W\cdot a(\mu)}{2}
  = \frac{R(1-\mu) - WR\mu - (1-W)\kappa^{\min}(\mu)}{2}
  \label{eq:transfer_mu}
\end{equation}
which is decreasing in $\mu$ when $W > 1$: higher equality redistributes
the expected lottery gain from $L$ to $H$, because $L$'s larger stake reduces
their net receipt. The risk premia respond to $\mu$ through the prize $\Pi(\mu)$
and the deviation $D_s(\mu)$:
\begin{align}
  \tilde{\rho}_L(\mu) &= \frac{\Pi(\mu)^2}{8} \cdot \frac{\alpha_u - 2\beta_u D_L(\mu)}{u'(\bar{b})}
  \label{eq:rhoL_mu} \\[4pt]
  \tilde{\rho}_H(\mu) &= \frac{\Pi(\mu)^2}{8} \cdot \frac{\alpha_u - 2\beta_u D_H(\mu)}{u'(\bar{b})}
  \label{eq:rhoH_mu}
\end{align}
where $D_L(\mu) = \bar{b}(\mu) - m^* + \epsilon_L(\mu)$ and
$D_H(\mu) = \bar{b}(\mu) - m^* - \epsilon_L(\mu)$.
As $\mu$ rises, $D_L$ moves toward zero from below (worker $L$ approaches
the inflection point) so $|\tilde\rho_L|$ falls --- the risk-seeking premium
diminishes as worker $L$ becomes richer and more moderate in risk attitude.
Conversely $D_H$ moves away from zero, so $\tilde\rho_H$ rises --- worker $H$
becomes more risk-averse as their relative income advantage grows.

\begin{proposition}[Risk premium responses to inequality]
Under the linear approximation of $u''$:
\begin{enumerate}
  \item The magnitude of worker $L$'s risk premium $|\tilde\rho_L|$ is
        decreasing in $\mu$: greater equality reduces risk-seeking behaviour
        for the lower-income worker.
  \item Worker $H$'s risk premium $\tilde\rho_H$ is increasing in $\mu$:
        greater equality increases risk-aversion for the higher-income worker.
  \item The net social risk premium $\tilde\rho_L + \tilde\rho_H = -\Pi^2 A(\bar{b})/4$
        is positive (lottery socially desirable) iff $\bar{b}(\mu) < m^*$.
        Since $\bar{b}(\mu) = \bar{b}_0 R\mu$ is increasing in $\mu$, there
        exists a critical $\mu^\dagger = m^*/(\bar{b}_0 R)$ above which the
        lottery is no longer socially desirable.
\end{enumerate}
\end{proposition}

\subsection{Summary: Worker Responses to $\mu$}

Table~\ref{tab:worker_responses} collects the signed responses of all
worker-side equilibrium objects to a marginal increase in wage equality $\mu$.

\begin{table}[h!]
\centering
\caption{Worker responses to an increase in wage equality ($\mu \uparrow$)}
\label{tab:worker_responses}
\begin{tabular}{llll}
\toprule
Object & Symbol & Response to $\mu\uparrow$ & Mechanism \\
\midrule
Worker $L$ income & $M_L = R\mu$ & $+$ & Direct \\
Worker $H$ income & $M_H = R(1-\mu)$ & $-$ & Direct \\
Upskilling floor & $\kappa^{\min}(\mu)$ & $+$ & Rising subsistence $\bar{b}_0 R\mu$ \\
Worker $L$ upskill frac. & $\lambda_L^*$ & $+$ (cond.) & See (\ref{eq:lL_dmu_condition}) \\
Worker $H$ upskill frac. & $\lambda_H^*$ & $+$ (always) & Both effects same sign \\
Total upskilling & $K^* = 2\kappa^{\min}$ & $+$ & Rising floor \\
Worker $L$ lottery stake & $a = M_L - \kappa^{\min}$ & $+$ & Income rises faster than floor \\
Worker $H$ lottery stake & $c = M_H - \kappa^{\min}$ & $-$ & Income falls and floor rises \\
Lottery pool & $\Lambda^* = Wa + c$ & $+$ & Income effect dominates \\
Net transfer to $L$ & $\epsilon_L^*$ & $-$ & Larger $L$ stake hurts $L$ \\
Worker $L$ risk premium & $|\tilde\rho_L|$ & $-$ & $L$ closer to inflection \\
Worker $H$ risk premium & $\tilde\rho_H$ & $+$ & $H$ further from inflection \\
Social premium & $\tilde\rho_L+\tilde\rho_H$ & Changes sign at $\mu^\dagger$ & $\bar{b}$ crosses $m^*$ \\
\bottomrule
\end{tabular}
\end{table}

%----------------------------------------------------------------------
\section{Question 2: What Constraints Does $\mu$ Place on Long-Term Growth?}
%----------------------------------------------------------------------

\subsection{The Firm's Asset Accumulation Problem}

The firm accumulates assets $A_t$ according to:
\begin{equation}
  A_{t+1} = \left[(1-\pi-\delta_f) + \pi\Phi(\mu_t)\right]A_t
             - K(\mu_t) - C(A_t)
  \label{eq:firm_lom}
\end{equation}
where:
\begin{align}
  \Phi(\mu_t) &= \frac{(W+1)\kappa^{\min}(\mu_t)}{R}
  \quad \text{(upskilling index, increasing in $\mu_t$)}
  \label{eq:Phi} \\[4pt]
  K(\mu_t) &= k|\mu_t - \mu_0|^\eta
  \quad \text{(restructuring cost, $k,\eta>0$)}
  \label{eq:K} \\[4pt]
  C(A_t) &= c_f A_t^\xi
  \quad \text{(maintenance cost, $c_f,\xi>0$)}
  \label{eq:C}
\end{align}
The parameter $\mu_0$ is the reference market wage split. $K(\mu_t)$ is
zero when the firm sets wages at the market reference and increases convexly
as the firm deviates from it. $\Phi(\mu_t)$ captures how the aggregate
upskilling return --- weighted by the population mass ratio $(W+1)$ ---
feeds into asset growth through the wage bill $\pi A_t$.

\subsection{The Upskilling Index and the Growth Mechanism}

The upskilling index $\Phi(\mu) = (W+1)\kappa^{\min}(\mu)/R$ encodes the
central link from wage inequality to growth. Substituting
$\kappa^{\min}(\mu)$ from (\ref{eq:kappa_min_gamma2}):
\begin{equation}
  \Phi(\mu) = \frac{(W+1)}{R} \cdot
  \frac{-1+\sqrt{1+4(\bar{b}_0 R\mu+\delta_w)}}{2}
  \label{eq:Phi_explicit}
\end{equation}
This is strictly increasing and strictly concave in $\mu$:
\begin{equation}
  \frac{d\Phi}{d\mu} = \frac{(W+1)\Psi(\mu)}{R} > 0,
  \qquad
  \frac{d^2\Phi}{d\mu^2} = \frac{(W+1)}{R}\frac{d\Psi}{d\mu} < 0
  \label{eq:Phi_derivs}
\end{equation}
The concavity of $\Phi$ in $\mu$ is the fundamental tension in the growth
problem. Greater equality raises upskilling returns (positive for growth)
but at a diminishing rate, while the restructuring cost $K(\mu)$ rises
convexly. The optimal $\mu$ trades off these two effects.

\subsection{The Stackelberg Timing and Non-Circularity}

The firm solves a Stackelberg problem: it moves first by setting $\mu_t$,
internalising how workers will best-respond with $\lambda_s^*(\mu_t)$.
By substituting the workers' best responses into $\Phi(\mu_t)$ before
optimising, the firm's problem reduces to a single-variable maximisation
with no circularity:
\begin{equation}
  \mu_t^* = \arg\max_{\mu_t} \left\{
    \left[(1-\pi-\delta_f) + \frac{\pi(W+1)\kappa^{\min}(\mu_t)}{R}\right]A_t
    - k|\mu_t-\mu_0|^\eta - c_f A_t^\xi
  \right\}
  \label{eq:firm_opt}
\end{equation}

\subsection{The Firm's First-Order Condition}

Differentiating (\ref{eq:firm_opt}) with respect to $\mu_t$ and setting to
zero, using $\eta = 2$ for closed form:
\begin{equation}
  \underbrace{\frac{\pi(W+1)\Psi(\mu_t^*)}{R} \cdot A_t}_{\text{marginal upskilling gain}}
  = \underbrace{2k(\mu_t^* - \mu_0)}_{\text{marginal restructuring cost}}
  \label{eq:foc}
\end{equation}
Solving for $\mu_t^*$, and evaluating $\Psi$ at the reference point
$\Psi_0 \equiv \Psi(\mu_0)$ for tractability:
\begin{equation}
  \boxed{\mu_t^* = \mu_0 + \Gamma A_t,
  \qquad \Gamma \equiv \frac{\pi(W+1)\Psi_0}{2kR}}
  \label{eq:mu_opt}
\end{equation}
This is the central result of the firm's problem. The optimal wage equality
$\mu_t^*$ is linear in the asset stock $A_t$. Richer firms set greater
equality. The coefficient $\Gamma$ is the key growth-inequality link:
it is increasing in the wage share $\pi$, the population mass ratio $W$,
and the subsistence sensitivity $\Psi_0$, and decreasing in the
restructuring cost $k$.

\begin{proposition}[Optimal wage inequality and firm wealth]
The firm's optimal income split $\mu_t^*$ is:
\begin{enumerate}
  \item Increasing in $A_t$: wealthier firms choose greater equality.
  \item Increasing in $W$: a larger lower-rank population pushes the firm
        toward equality.
  \item Increasing in $\pi$: a larger wage share raises the marginal growth
        benefit of equality.
  \item Decreasing in $k$: higher restructuring costs push the firm toward
        the market reference $\mu_0$.
  \item Always above $\mu_0$ when $A_t > 0$: the firm always chooses
        greater equality than the market reference.
\end{enumerate}
\end{proposition}

\subsection{The Constraints on $\mu$ and the Viability Condition}

Not every value of $\mu$ supports a positive steady-state asset level.
Setting $A_{t+1} = A_t = A^*$ in (\ref{eq:firm_lom}) and using
$\mu^* = \mu_0 + \Gamma A^*$:
\begin{equation}
  A^*[\pi\Phi^* - \delta_f - c_f] = k(\mu^* - \mu_0)^2 = k\Gamma^2 A^{*2}
  \label{eq:ss_eq}
\end{equation}
Dividing by $A^* > 0$:
\begin{equation}
  \pi\Phi^* - \delta_f - c_f = k\Gamma^2 A^*
  \label{eq:ss_div}
\end{equation}
For a positive $A^*$ to exist, the left-hand side must be positive. Evaluating
$\Phi^*$ at the reference $\kappa_0 = \kappa^{\min}(\mu_0)$ gives the
\emph{viability condition}:
\begin{equation}
  \boxed{\pi(W+1)\kappa_0 > R(\delta_f + c_f)}
  \label{eq:viability}
\end{equation}
This is the fundamental constraint that $\mu$ --- through $\kappa_0 = \kappa^{\min}(\mu_0)$
--- places on long-term growth. Condition (\ref{eq:viability}) says: the
upskilling return, amplified by the wage share and population mass, must
exceed the total running costs of the firm. When it fails, no amount of
wage restructuring can save the firm --- the growth mechanism is simply
too weak. When it holds, the steady-state asset level is:
\begin{equation}
  \boxed{A^* = \frac{\pi(W+1)\kappa_0/R - \delta_f - c_f}{k\Gamma^2}}
  \label{eq:Astar}
\end{equation}

\subsection{The Steady-State Equilibrium}

Given $A^*$, the full steady-state equilibrium is determined sequentially
with no implicit equations remaining (for $\gamma=2$, $\eta=2$):

\begin{align}
  \text{Derived constants:}\quad
  & \kappa_0 = \frac{-1+\sqrt{1+4(\bar{b}_0 R\mu_0+\delta_w)}}{2}
  \label{eq:k0} \\[4pt]
  & \Psi_0 = \frac{\bar{b}_0 R}{1+2\kappa_0}
  \label{eq:Psi0} \\[4pt]
  & \Gamma = \frac{\pi(W+1)\Psi_0}{2kR}
  \label{eq:Gamma} \\[12pt]
  \text{Firm steady state:}\quad
  & A^* = \frac{\pi(W+1)\kappa_0/R - \delta_f - c_f}{k\Gamma^2}
  \label{eq:Astar2} \\[4pt]
  & \mu^* = \mu_0 + \Gamma A^*
  \label{eq:mustar} \\[12pt]
  \text{Worker steady state:}\quad
  & \kappa^* = \frac{-1+\sqrt{1+4(\bar{b}_0 R\mu^*+\delta_w)}}{2}
  \label{eq:kstar} \\[4pt]
  & \lambda_L^* = \frac{\kappa^*}{R\mu^*},
  \quad \lambda_H^* = \frac{\kappa^*}{R(1-\mu^*)}
  \label{eq:lambdas} \\[4pt]
  & K^* = 2\kappa^*, \quad
  a^* = R\mu^*-\kappa^*, \quad
  c^* = R(1-\mu^*)-\kappa^*
  \label{eq:Kac} \\[4pt]
  & \Lambda^* = Wa^*+c^*, \quad
  \epsilon_L^* = \tfrac{c^*-Wa^*}{2}
  \label{eq:Lam_eL} \\[4pt]
  \text{Risk premia:}\quad
  & \tilde\rho_s^* = \frac{\Lambda^{*2}}{8}\cdot\frac{\alpha_u - 2\beta_u D_s^*}{u'(\bar{b}^*)}
  \label{eq:rhostar}
\end{align}

Every object is an explicit function of the nine primitives
$(R, \mu_0, W, \bar{b}_0, \delta_w, \pi, \delta_f, k, c_f)$
plus the utility parameters $(\alpha_u, \beta_u, m^*)$.

\subsection{The $\mu$-Growth Relationship and its Constraints}

Combining the results, the asset growth rate is:
\begin{equation}
  g(\mu, A) = \pi\Phi(\mu) - \delta_f - c_f - \frac{k(\mu-\mu_0)^2}{A}
  \label{eq:growth_rate}
\end{equation}
This expression characterises precisely how $\mu$ constrains growth:

\textbf{Lower bound on $\mu$.} The growth rate is increasing in $\mu$ for
$\mu < \mu^*(A)$, since $\Phi$ is increasing and the restructuring cost is
below its maximum. Setting $g = 0$ at $\mu_0$ (the minimum feasible $\mu$
with no restructuring cost) gives the viability condition (\ref{eq:viability}).
If this fails even at $\mu = \mu_0$, no positive growth is achievable and
the firm collapses.

\textbf{Upper bound on $\mu$.} Feasibility requires $\mu < \tfrac{1}{2}$
(worker $H$ must earn more than worker $L$). Additionally, the restructuring
cost $k(\mu-\mu_0)^2$ must not exceed the upskilling gain
$\pi[\Phi(\mu)-\Phi(\mu_0)]A$. This gives an effective upper bound:
\begin{equation}
  \mu \leq \mu_0 + \sqrt{\frac{\pi[\Phi(\mu)-\Phi(\mu_0)]A}{k}}
  \label{eq:upper_bound}
\end{equation}
which is increasing in $A$ --- richer firms can sustain more redistribution.

\textbf{The growth-maximising $\mu$.} Differentiating $g$ with respect to $\mu$
and setting to zero recovers exactly $\mu_t^* = \mu_0 + \Gamma A$ from
(\ref{eq:mu_opt}). Growth maximisation and profit maximisation coincide.
Since $\partial^2 g/\partial\mu^2 < 0$ (concave $\Phi$ plus convex cost),
the growth-maximising $\mu^*$ is unique and interior.

\textbf{Endogenous Kuznets relationship.} At stationarity,
$\mu^* = \mu_0 + \Gamma A^*$, so wealthier economies (higher $A^*$) operate
at a higher $\mu^*$ (greater equality). Along the transitional path,
as $A_t$ grows from $A_0$ toward $A^*$, $\mu_t^*$ rises monotonically.
The Gini coefficient of income (proxied by $1-2\mu^*$) falls as the
economy grows. This endogenous Kuznets relationship --- inequality falls
as income rises --- is a direct consequence of the $\mu$-growth mechanism.

%----------------------------------------------------------------------
\section{Full Comparative Statics}
%----------------------------------------------------------------------

Table~\ref{tab:cs} presents the complete comparative statics of the
steady-state equilibrium, derived analytically from
equations (\ref{eq:k0})--(\ref{eq:rhostar}).

\begin{table}[h!]
\centering
\caption{Steady-state comparative statics --- all three players}
\label{tab:cs}
\begin{tabular}{lcccccccc}
\toprule
Shock & $\partial A^*/\partial x$ & $\partial\mu^*/\partial x$
      & $\partial\lambda_L^*/\partial x$ & $\partial\lambda_H^*/\partial x$
      & $\partial K^*/\partial x$ & $\partial\Lambda^*/\partial x$
      & $\partial\epsilon_L^*/\partial x$ & $\partial(\tilde\rho_L^*+\tilde\rho_H^*)/\partial x$ \\
\midrule
$W\uparrow$      & $+$ & $+$ & amb. & $-$   & amb. & $+$   & $-$   & $+$ \\
$\pi\uparrow$    & $+$ & $+$ & $+$  & $+$   & $+$  & $+$   & amb.  & $+$ \\
$k\uparrow$      & $-$ & $-$ & $-$  & $-$   & $-$  & $-$   & amb.  & $-$ \\
$\bar{b}_0\uparrow$ & amb. & $+$ & $+$ & $-$  & $+$  & amb.  & $-$   & $-$ \\
$\delta_f\uparrow$ & $-$ & $-$ & $-$ & $-$   & $-$  & $-$   & $+$   & $-$ \\
$\delta_w\uparrow$ & $+$ & $+$ & $+$ & amb.  & $+$  & $+$   & $-$   & $+$ \\
$\mu_0\uparrow$  & $+$ & $+$ & amb. & $-$   & $+$  & $+$   & $-$   & $+$ \\
$R\uparrow$      & amb. & amb. & $-$ & $-$   & amb. & amb.  & amb.  & amb. \\
$\alpha_u\uparrow$ & --- & --- & --- & ---   & ---  & ---   & ---   & $-$ \\
$\beta_u\uparrow$  & --- & --- & --- & ---   & ---  & ---   & ---   & $+$ \\
$m^*\uparrow$      & --- & --- & --- & ---   & ---  & ---   & ---   & $+$ \\
\midrule
\end{tabular}
\medskip

\small
Notes: ``amb.'' = ambiguous sign, parameter-dependent.
Dashes (---) indicate the parameter does not enter the firm or
worker allocation problem directly, only the risk premium through utility curvature.
$\tilde\rho_L^*+\tilde\rho_H^* = -\Lambda^{*2}A(\bar{b}^*)/4$;
its sign responses reflect changes in both $\Lambda^*$ and $A(\bar{b}^*)$.
\end{table}

Several patterns in Table~\ref{tab:cs} are worth highlighting.

\textbf{Wage share $\pi$.}
A higher wage share unambiguously raises all steady-state quantities except
possibly $\epsilon_L^*$. More of the asset base is converted into wages,
which raises the upskilling return and drives the firm toward greater equality.
This is the strongest positive lever on growth in the model.

\textbf{Restructuring cost $k$.}
Higher restructuring costs reduce everything. They discourage the firm from
deviating from $\mu_0$, reducing the upskilling return and compressing $A^*$.
This is the primary friction limiting the firm's ability to use wage policy
for growth.

\textbf{Population mass $W$.}
Greater lower-rank population mass raises $A^*$ and $\mu^*$ --- the firm
responds to a larger lower-rank group by raising equality, which raises the
upskilling index. Worker $H$'s upskilling fraction falls because their
effective income $R(1-\mu^*)$ falls as $\mu^*$ rises. The lottery pool
$\Lambda^*$ rises because $W$ directly amplifies $L$'s contribution.

\textbf{Utility parameters $\alpha_u$, $\beta_u$, $m^*$.}
These affect only the risk premia, not the allocation or accumulation objects.
Greater curvature $\alpha_u$ reduces the net social premium (more risk aversion
overall), while greater skewness $\beta_u$ raises it (more pronounced
Friedman-Savage asymmetry). A higher inflection $m^*$ raises the premium
because it pushes both workers into regions of stronger curvature. This
separation of utility parameters from accumulation outcomes is a direct
consequence of the linear approximation: risk preferences govern the
premium but not the equilibrium quantities.

%----------------------------------------------------------------------
\section{Conclusion}
%----------------------------------------------------------------------

We have derived a tractable three-player model in which wage inequality $\mu$
is the central variable connecting worker behaviour to firm growth. The model
delivers closed-form answers to both research questions.

On the first question --- how workers respond to $\mu$ --- the main findings
are: (i) both workers' upskilling fractions respond positively to greater
equality under empirically relevant conditions; (ii) total upskilling $K^*$
is increasing in $\mu$ through the rising subsistence floor, but the aggregate
neutral to pure redistribution holding the floor fixed; (iii) worker $L$'s
risk premium falls in magnitude as $\mu$ rises (approaching the inflection),
while worker $H$'s risk premium rises; (iv) the lottery is socially desirable
if and only if subsistence is below the Friedman-Savage inflection point,
giving a critical threshold $\mu^\dagger = m^*/(\bar{b}_0 R)$.

On the second question --- what constraints $\mu$ places on growth --- the
main findings are: (i) the viability condition $\pi(W+1)\kappa_0 > R(\delta_f+c_f)$
is the hard floor on growth: if it fails, no wage policy can generate a
positive steady state; (ii) conditional on viability, the firm's optimal
$\mu^*$ is linear in the asset stock, so richer firms enforce greater
equality --- an endogenous Kuznets relationship; (iii) the growth-maximising
and profit-maximising $\mu^*$ coincide, so there is no agency conflict
between firm goals and worker welfare on this margin; (iv) a poverty trap
exists for low initial asset levels, distinguishing the model from
frameworks with unconditional convergence.

The linear approximation of $u''$ is the methodological contribution that
makes all of this tractable. It preserves the sign asymmetry between workers'
risk premia --- the essential Friedman-Savage content --- while keeping
equilibrium utilities, risk premia, and comparative statics fully explicit.
The approximation introduces error only in the magnitude of risk premia, not
in their signs or qualitative comparative statics, and the error is small
when lottery prizes are modest relative to the distance from subsistence to
the inflection point.

%----------------------------------------------------------------------
\appendix
\section{Complete Notation Table}
%----------------------------------------------------------------------

\begin{table}[h!]
\centering
\caption{All symbols, definitions, and domains}
\label{tab:notation}
\begin{tabular}{lp{8.5cm}l}
\toprule
Symbol & Definition & Domain \\
\midrule
\multicolumn{3}{l}{\textit{Primitives}} \\
$R$ & Total income available to workers & $\mathbb{R}_{++}$ \\
$\mu$ & Income share of worker $L$; wage inequality parameter & $(0,\tfrac{1}{2})$ \\
$\mu_0$ & Reference (market) income split & $(0,\tfrac{1}{2})$ \\
$W$ & Population mass ratio $= i^*/(N-i^*)$ & $[1,\infty)$ \\
$\bar{b}_0$ & Subsistence coefficient; $\bar{b} = \bar{b}_0 R\mu$ & $\mathbb{R}_{++}$ \\
$\delta_w$ & Human capital depreciation & $[0,1)$ \\
$\gamma$ & Upskilling production power & $(0,1)$ \\
$\pi$ & Firm wage share of assets & $(0,1)$ \\
$\delta_f$ & Firm asset depreciation rate & $(0,1)$ \\
$k$ & Restructuring cost scale & $\mathbb{R}_{++}$ \\
$\eta$ & Restructuring cost power & $(1,\infty)$ \\
$c_f$ & Asset maintenance cost scale & $\mathbb{R}_{++}$ \\
$\xi$ & Asset maintenance cost power & $(1,\infty)$ \\
$\alpha_u$ & Friedman-Savage curvature & $\mathbb{R}_{++}$ \\
$\beta_u$ & Friedman-Savage skewness & $\mathbb{R}_{++}$ \\
$m^*$ & Friedman-Savage inflection income & $\mathbb{R}_{++}$ \\
\midrule
\multicolumn{3}{l}{\textit{Worker objects}} \\
$M_L, M_H$ & Worker incomes $= R\mu,\, R(1-\mu)$ & $\mathbb{R}_{++}$ \\
$\lambda_s$ & Fraction of income to upskilling & $[0,1]$ \\
$\kappa_s$ & Upskilling investment $= \lambda_s M_s$ & $\mathbb{R}_{+}$ \\
$\ell_s$ & Lottery stake $= (1-\lambda_s)M_s$ & $\mathbb{R}_{+}$ \\
$b_s$ & Base income $= \kappa_s+\kappa_s^\gamma-\delta_w$ & $\mathbb{R}$ \\
$\bar{b}$ & Subsistence floor $= \bar{b}_0 R\mu$ & $\mathbb{R}_{++}$ \\
$\kappa^{\min}(\mu)$ & Minimum upskilling; solves $\kappa+\kappa^\gamma = \bar{b}+\delta_w$ & $\mathbb{R}_{++}$ \\
$\lambda_s^*$ & Optimal upskilling fraction at constraint & $(0,1)$ \\
$a$ & Worker $L$ lottery stake $= R\mu-\kappa^{\min}$ & $\mathbb{R}_{++}$ \\
$c$ & Worker $H$ lottery stake $= R(1-\mu)-\kappa^{\min}$ & $\mathbb{R}_{++}$ \\
$\Lambda^*$ & Total lottery pool $= Wa+c$ & $\mathbb{R}_{++}$ \\
$\epsilon_L^*$ & Net expected transfer to $L$ $= (c-Wa)/2$ & $\mathbb{R}$ \\
$K^*$ & Total upskilling $= 2\kappa^{\min}$ & $\mathbb{R}_{++}$ \\
$D_s$ & Deviation $= \bar{b}-m^*+\epsilon_s$ & $\mathbb{R}$ \\
$\tilde\rho_s$ & Linearly approx.\ risk premium & $\mathbb{R}$ \\
$V_s$ & Worker expected utility & $\mathbb{R}$ \\
\midrule
\multicolumn{3}{l}{\textit{Firm objects}} \\
$A_t$ & Asset stock at period $t$ & $\mathbb{R}_{+}$ \\
$\Phi(\mu)$ & Upskilling index $= (W+1)\kappa^{\min}(\mu)/R$ & $\mathbb{R}_{++}$ \\
$\Psi(\mu)$ & Upskilling sensitivity $= d\kappa^{\min}/d\mu$ & $\mathbb{R}_{++}$ \\
$\Psi_0$ & $\Psi$ evaluated at $\mu_0$: $= \bar{b}_0 R/(1+2\kappa_0)$ & $\mathbb{R}_{++}$ \\
$\kappa_0$ & $\kappa^{\min}$ evaluated at $\mu_0$ & $\mathbb{R}_{++}$ \\
$\Gamma$ & Firm FOC coefficient $= \pi(W+1)\Psi_0/(2kR)$ & $\mathbb{R}_{++}$ \\
$\mu_t^*$ & Optimal wage split at time $t$ $= \mu_0+\Gamma A_t$ & $(0,\tfrac{1}{2})$ \\
$A^*$ & Steady-state asset level & $\mathbb{R}_{+}$ \\
$\mu^*$ & Steady-state wage split & $(0,\tfrac{1}{2})$ \\
$\mu^\dagger$ & Social premium threshold $= m^*/(\bar{b}_0 R)$ & $(0,\tfrac{1}{2})$ \\
\bottomrule
\end{tabular}
\end{table}

\end{document}
